%% kapa: exam=1 round=35 year=2024 period=2 form=A subject=accounting subject_ko=회계학 seq=30 number=70
%% source: Q-Net 감정평가사 기출문제 (cst00309) artlSeq=5213229 file=raw/1008/5213229/2교시 문제지 원본.pdf
%% conversion: kordoc 4.24.51 -> tools/gen1.py (auto); review=proofread-sonnet; flags=tab
\begin{question}{70}
(주)감평과 (주)한국은 사용 중인 유형자산을 상호 교환하여 취득하였다. 동 교환 거래에서 (주)한국의 유형자산 공정가치가 (주)감평의 유형자산 공정가치보다 더 명백하며, (주)감평은 (주)한국으로부터 추가로 현금 ￦3,000을 수취하였다. 두 회사가 보유하고 있는 유형자산의 장부금액과 공정가치가 다음과 같을 때, (주)감평과 (주)한국이 인식할 유형자산처분손익은? (단, 두 자산의 공정가치는 신뢰성 있게 측정할 수 있으며, 상업적 실질이 있다.)
\par{\small\centering\begin{tabular}{|c|c|c|}\hline
구분 & (주)감평 & (주)한국 \tabularnewline\hline
장부금액(순액) & ￦10,000 & ￦8,000 \tabularnewline
공정가치 & 9,800 & 7,900 \tabularnewline\hline
\end{tabular}\par}
\smallskip
\noindent\begin{tabular}{@{}l@{\hspace{0.4em}}p{0.17\linewidth}p{0.17\linewidth}@{\hspace{1em}}l@{\hspace{0.4em}}p{0.17\linewidth}p{0.17\linewidth}@{}}
 & \uline{(주)감평} & \uline{(주)한국} & & \uline{(주)감평} & \uline{(주)한국} \tabularnewline
① & 손실 ￦200 & 손실 ￦100 & ② & 손실 ￦200 & 손실 ￦1,200 \tabularnewline
③ & 이익 ￦200 & 이익 ￦900 & ④ & 이익 ￦900 & 손실 ￦100 \tabularnewline
⑤ & 이익 ￦900 & 손실 ￦1,200 & & & \tabularnewline
\end{tabular}
\end{question}
\answer{4}
